\documentclass[tikz, border=5mm]{article}

\usepackage{geometry}
\usepackage{amsmath}
\usepackage{graphicx}
\usepackage{listings}%For code do \begin{lstlisting}
\usepackage{color}
\usepackage{tikz-cd}
\usepackage{tikz}
\usepackage{float}
\usepackage{glossaries}
\usepackage{fancyhdr}
\usepackage{cleveref}
\usepackage{epigraph}
\usetikzlibrary{arrows,arrows.meta,backgrounds,shapes,automata,petri,positioning,calc,chains,decorations.pathmorphing,fit,decorations.pathreplacing}
\graphicspath{ {../assets/} }

\definecolor{dkgreen}{rgb}{0,0.6,0}
\definecolor{gray}{rgb}{0.5,0.5,0.5}
\definecolor{mauve}{rgb}{0.58,0,0.82}
\definecolor{backcolour}{rgb}{0.95,0.95,0.92}

\DeclareRobustCommand{\comment}[1]{
    {\color{orange}\emph{#1}}
}

\newcounter{address}
\newcommand\StateDiagram[4]{
\begin{tikzpicture}[>=Stealth, start chain=a, start chain=b, node distance=0,
  box/.style={draw, on chain=##1, minimum width=15mm, minimum height=6mm},
  label/.style={thick,minimum width=2cm,fill=blue!20}]
  \begin{scope}[local bounding box=store]
  \node [box=a,label] (a) {Address};
  \node [box=a] {$\cdots$};

  \ifnum#1>99\relax\else
  \foreach \x in {#1,...,99}
  {
    \node [box=a] (addr_\x) {\x};
  }
  \fi
  \node [box=b,below=of a,label] {Value} ;
  \node [box=b] {$\cdots$} ;
  \setcounter{address}{#1}
  \foreach \x in {#3}
  {
    \node [box=b] (mem_\theaddress) {\x};
    \stepcounter{address}
  }
  \node [above=of addr_#2,yshift=5mm] (freelabel) { free } ;
  \draw [->] (freelabel) -- (addr_#2.north) ;
  \end{scope}
  \node [left=of a, anchor=north east] {Store};
  \begin{scope}[start chain=c, start chain=d,shift={($(store.south)-(0,1cm)$)}]
    \node [box=c,label] (c) {Variable} ;
    % varname, addr
    \foreach \x/\y in {#4}
    {
      \node [box=c] (var_\x) {\x};
    }
    \node [box=d,below=of c,label] {Address} ;
    \foreach \x/\y in {#4}
    {
      \node [box=d] {\y};
    }
  \end{scope}
  \node [left=of c, anchor=north east] {Environment};
\end{tikzpicture}
}
{
\newglossaryentry{DSL}{name={\textbf{Domain Specific Languages}},description={
        A language(i.e not just a library) with abstractions targetet at a specific problem domain.
        \begin{itemize}
            \item \textit{External DSL} - A DSL defined as a seperate programming language.
            \item \textit{Internal/Embedded DSL} - A DSL defined as a language-like interface to a library.
        \end{itemize} 
    }}
\newglossaryentry{GPL}{name={\textbf{General Purpose Languages}},description={
    A language suited for a wide variety of problems and situations but lacks specialized features to deal with specific programs like a DSL.
}}
\newglossaryentry{syntax}{name={\textbf{Syntax}},description={Syntax refers to the rules that define the structure of a language. 
                                                    Syntax in computer programming means the rules that control the structure of the symbols, 
                                                    punctuation, and words of a programming language.}}

\newglossaryentry{semantics}{name={\textbf{Semantics}},description={The semantics of a program concerns the meaning of the program. i.e what it actually does. 
                                                            It can take many forms, sometimes we're only interested in the result of the program.
                                                            Other cases concern the steps taken by the program to reach said output.}}
\newglossaryentry{interpreter}{name={\textbf{Interpreter}},description={An interpreter is a program that directly executes instructions written in a programming language.}}
\newglossaryentry{compiler}{name={\textbf{Compiler}},description={A compiler is a program that translates computer code(\textit{the source language}) into another language(\textit{the target language})
                                                                Compilers usually convert some high level language to some lower level language.}}
\newglossaryentry{AST}{name={\textbf{Abstract Syntax Tree}},description={An Abstract Syntax Tree(AST) is a tree representation of the syntactic structure of our program. 
                                                                Can be represented by using trees or terms, and described by an algebraic data type or regular tree grammar.}}
\newglossaryentry{meta-programming}{name={\textbf{Meta-Programming}},description={Metaprogramming is a programming technique in which computer programs have the ability to treat other programs as their data. 
                                        It means that a program can be designed to read, generate, analyze or transform other programs, and even modify itself while running.}}
\newglossaryentry{BSL}{name={\textbf{Basic Signature Language}},description={Simple language to illustrate signatures}}
\newglossaryentry{BIPL}{name={\textbf{Basic Imperative Programing Language}},description={BIPL is a trivial language to explain basic imperative programming concepts like mutable vars, assignments, control-flow, loops, and iteration}}
\newglossaryentry{expression}{name={\textbf{Expression}},description={An expression is a syntactic construct that can be evaluated in order to obtain its value. The resulting value is usually one of the program's types.}}
\newglossaryentry{wellformed}{name={\textbf{Wellformed}},description={Wellformedness is when a program is following all the rules. smileemoji}}
\newglossaryentry{store}{name={\textbf{Store}},description={Program memory, byte/value array, grows uppwards.}}
\newglossaryentry{env}{name={\textbf{Enviroment}},description={Map describing where things are located in the \gls{store}. Kinda like a phonebook}}
\newglossaryentry{scope}{name={\textbf{Scope}},description={A collection of identifier bindings . i,e what's captured by the eniroment at some point in the code.}}
\newglossaryentry{param}{name={\textbf{Parameter}},description={A parameter is a localvariable that is initialized with the arguments. Also often contains how those args are to be treated.}}
\newglossaryentry{local var}{name={\textbf{Local Variable}},description={A local variable is a variable that only exists within a limited scope.}}
\newglossaryentry{stackframe}{name={\textbf{Stackframe}},description={The stackframe is a snapshot of how the program enviroment looks at a certain point in time. 
                                                                        The stackframe saves things that could be changed by running the procedure and lets us restore the program
                                                                        to the previous state without all the changes made by the procedure.}}
\newglossaryentry{argument}{name={\textbf{Argument}},description={An argument is a value provided to the procedure when it is run. When the procedure is run the parameters of the procedure is initialized with its corresponding argument.}}
\newglossaryentry{copy sem}{name={\textbf{Copy Semantics}},description={Type of argument passing where the parameters are initialized with the value of the arguments}}
\newglossaryentry{ref sem}{name={\textbf{Reference Semantics}},description={Type of argument passing where the parameters point to the the address as the argument.}}
\newglossaryentry{SLE}{name={\textbf{Software Language Engineering}},description={Software Language Engineering is the scientific field that researches language development, and mainteinace of formal descriptios, and tooling of software lanugages}}
\newglossaryentry{BNF}{name={\textbf{Backus-Naur Form}},description={Formal notation for describing grammars. Used to describe the syntax of a language.}}
\newglossaryentry{sum of prod}{name={\textbf{Sum of Products}},description={}}
}
\makenoidxglossaries
%For handling code, addapt depending on language
\lstset{frame=tb,
  language=PASCAL,
  aboveskip=3mm,
  belowskip=3mm,
  showstringspaces=false,
  columns=flexible,
  basicstyle={\small\ttfamily},
  numbers=left,
  backgroundcolor=\color{backcolour}, 
  numberstyle=\tiny\color{gray},
  keywordstyle=\color{blue},
  commentstyle=\color{dkgreen},
  stringstyle=\color{mauve},
  breaklines=true,
  breakatwhitespace=true,
  tabsize=3
}
\pagestyle{fancy}
\fancyhf{}
\rhead{\rightmark}
\lhead{\leftmark}
\rfoot{\thepage}

\begin{document}
    \begin{titlepage}
        \begin{center}
            \vspace*{1cm}
    
            \huge
            \textbf{Book of Magne}
    
            \vspace{0.5cm}
            \LARGE
            INF222 Crashcourse 2022
                
            \vspace{1.5cm}

            \textbf{Sander Wiig}
    
            \vfill
            
            \Large
            INF222 Crashcourse for v2022.\\
            Some sections are adapted from Anya's INF225 notes and course material.
                
            \vspace{0.5cm}
        
            \includegraphics[width=0.4\textwidth]{UiBlogoMN_gray_m_Eng.png}\\
            \Large
            Institute for Informatics\\
            University of Bergen\\
            Norway\\
            \today
                
        \end{center}
    \end{titlepage}
    \newpage
    \tableofcontents
    \newpage
    \section{Introduction}
    \epigraph{We're building something here, detective. We're building it from scratch. All the pieces matter.}{\textit{Lester Freemont \\ The Wire}}
    \comment{TODO: \#1 Write introduction}
    \section{The Basics}
    \subsection{What is a programming language}
        A programming language 
        \begin{itemize}
            \item is an artifical language(i.e made by us humans on purpose)
            \item used to tell machines what to do
        \end{itemize}
        More formally a programming language is a set of rules that converts some input, like strings, into instructions that the computer can follow.
        This is of course a very general description and it, therefore, follows that there are many different types of programming languages.
        IT therefore should come as little surprise that we group languages by features and properties. 
        \subsubsection{Types of languages}
        There are many ways of grouping languages. They can be grouped by Purpose, typing, paradigm, Generality vs. Specificity, and many more.
        For now, we're going to group them by paradigm, and Generality vs. Specificity.
        \subsubsection*{Generality vs. Specificity}
        Languages are usually grouped into two categories when based on their specificity.
        \begin{itemize}
            \item DSL
            \item GPL
        \end{itemize}

        \Gls{DSL} are as the name suggests languages with a "specific" domain.\\
        DSLs usually have limited scope and use. Examples are JSON and SQL.     
        A Domain-Specific Language is a programming language with a higher level of abstraction optimized for a specific class of problems. 
        Optimized for a certain problem/domain. 
        DSLs can be further subdivided into external DSL(separate programming languages), and internal DLS(language-like interface as a library.)

        \Gls{GPL} however are more general and can be used to solve many different problems in many different situations.
        These languages have a wide array
        of uses and are usually what we think of when we hear the words programming language. Examples of GPLs are Java and Haskell.\\
        \begin{figure*}[!h]
            \begin{tabular}{|c|c|c}
                \hline
                \textbf{Characteristic} &\textbf{DSL vs. GPL}\\
                \hline
                \textbf{Domain} & DSLs have a small and well-defined domain\\
                \hline
                \textbf{Size} &GPLs are large, DSLs are usually small\\
                \hline
                \textbf{Lifespan} &GPLs last for years to decades, DSLs typicaly live for shorter periors.\\
                \hline
            \end{tabular}%
            \caption{Some more comparisons between GPLs and DSLs}
        \end{figure*}%

        \subsubsection*{Paradigm}
        We can also classify languages by programming paradigm, some of these are; 
        \begin{itemize}
            \item Imperative Languages, i.e C
            \item Functional Languages, i.e Haskell
            \item Object-oriented Languages, i.e Java, C\#, C++
            \item Logic Languages
        \end{itemize}
        \comment{Write something about this also}
        \comment{Write body}

        \subsubsection*{Syntax and Semantics}
        All programming languages have two parts; the \gls{syntax}, and the \gls{semantics}.\\
        Syntax is the study of \textit{structure}, just as semantics is the study of \textit{meaning}. Or in other words
        the syntax tells us \textit{how} to write legal programs, the semantics tells us \textit{what} those programs do. 

        \subsubsection{The Interpreter process}
            
            An \gls{interpreter} is a computer program that directly executes instructions written in a programming language.
            This differs from \glspl{compiler} since they don't have to translate the instructions into machine code to run it. 
            We use interpreters to define the semantics of our \gls{AST}(AST). They are relatively straightforward to implement since
            since each construct of the AST defines the semantics of that construct.\\
            \newline
            An interpreter is composed of several phases. \\
            The first is the lexical analyzer. The lexical analyzer takes the actual characters that the code is made up of and divides it up into its lexical tokens(by the tokenizer) using the concrete syntax of the program\footnote{Not covered by this course, so you can safely ignore how this works.}. 
            It then sends that token stream to the next part of the interpreter. The static analyzer.\\
            The static analyzer builds a parse tree out of the tokens and the concrete syntax. This is then converted into an AST by the abstract syntax rules. This AST is then given to the next part, the Semantic Analyser.
            The AST is type-checked, checked for well-formedness, names are resolved, types are inferred, and much more is discussed in chapter 3. The new AST, now with added information, is then given to the evaluator so that it can be evaluated and produce a result.
            See figure 2. for a visual description of the steps.
            \begin{figure}[!h]
                \centering
                \begin{minipage}{.5\textwidth}
                    \centering
                    \fbox{\centering
                    \begin{tikzcd}
                        \parbox{\textwidth}{}\arrow[d, "\text{character stream}"]\\
                        \fbox{\parbox{10em}{\centering Lexical Analyzer}}\arrow[d, "\text{token stream}"]\\
                        \fbox{\parbox{10em}{\centering Syntax Analyzer}}\arrow[d, "\text{syntax tree(AST)}"]\\
                        \fbox{\parbox{10em}{\centering Semantic Analyzer}}\arrow[d,"\text{syntax tree(Annotated AST)}"]\\
                        \fbox{\parbox{10em}{\centering Evaluator}}\arrow[d, "\text{result}"]\\
                        \parbox{\textwidth}{}
                    \end{tikzcd}}
                    \caption{Phases of an interpreter}
                \end{minipage}%
                \begin{minipage}{.5\textwidth}
                    \centering
                    \fbox{\centering
                    \begin{tikzcd}
                        \parbox{\textwidth}{}\arrow[d, "\text{character stream}"]\\
                        \fbox{\parbox{10em}{\centering Lexical Analyzer}}\arrow[d, "\text{token stream}"]\\
                        \fbox{\parbox{10em}{\centering Syntax Analyzer}}\arrow[d, "\text{syntax tree(AST)}"]\\
                        \fbox{\parbox{10em}{\centering Semantic Analyzer}}\arrow[d,"\text{syntax tree(Annotated AST)}"]\\
                        \fbox{\parbox{10em}{\centering Intermediate Code Generator}}\arrow[d, "\text{Intermediate representation}"]\\
                        \fbox{\parbox{10em}{\centering Machine-Independent Code Optimization}}\arrow[d, "\text{Intermediate representation}"]\\
                        \fbox{\parbox{10em}{\centering Code Generator}}\arrow[d, "\text{target machine code}"]\\
                        \fbox{\parbox{10em}{\centering Machine-Dependent Code Optimization}}\arrow[d, "\text{target-machine code}"]\\
                        \parbox{\textwidth}{}
                    \end{tikzcd}}
                    \caption{Phases of a "normal" compiler}
                \end{minipage}%
                \caption{Please apprecitate the figures above, they were hard to make}
            \end{figure}%
        
    \subsection{Meta Programming}
    One of the harder things in the course is \gls{meta-programming}. INF222 is usually the first time you've encountered meta-programming and it can be hard a hard concept to grasp. 
    Meta-programming is programming \textit{about} programming. More properly meta-programms treat other programms as data. When you see a datastructure like \gls{BSL} or \gls{BIPL} in Haskell it
    represents a program.

    \subsection{Q\&A}

    \section{ASTs \& Static Analysis}
    \subsection{AST}
        An \gls{AST} is the tree representation of the syntactic structure of a program; The abstract syntax
        describes the structure of the abstract syntax tree - it can be defined using a
        regular tree grammar, or an algebraic data type or term (in Rascal, ML, Prolog,
        ...), or an object-oriented inheritance hierarchy of node classes (Java, C++, ...),
        or as an S-expression (in Lisp languages).
        The abstract syntax tree can be used as an internal representation in a language processor, 
        but it is not the only possible representation.\\
        An abstract syntax can be generated by a grammar in the following way:
        For every non-terminal type, there is a corresponding abstract syntax type.
        Each type has one constructor (or node type) corresponding to each production in the grammar, with one child for every symbol in the production that
        is not a literal token (e.g, punctuation, keywords, or spaces). If a constructor
        has only one child, of the same type, it can be removed (e.g., this would be
        the case for a parenthesis expression). You can do this process entirely based
        on the information contained in a parse tree. Translating a parse tree into a
        corresponding abstract syntax tree is called imploding the parse tree.
        Given an abstract syntax tree, it is possible to reconstruct a parse tree or program text, given the original grammar - though the resulting program may be
        slightly different in terms of spaces and punctuation. This is called unparsing or
        pretty-printing (particularly if the output is nicely formatted). Parsing, imploding, pretty-printing, and then reparsing may not yield the exact same parse tree
        as the original tree, but it should still implode to the same abstract syntax tree (otherwise there is a bug in your toolchain!).

        Various phases in a language processor may change the abstract syntax
        tree, or use slightly different versions of the abstract syntax (e.g., after type
        checking, the nodes for variables include the type of the variable) - it is also
        possible to decorate or annotate the AST as processing proceeds. 
        This adds extra information to the nodes in the AST, without impacting the structure of the abstract syntax.
        Important abstract syntax design considerations are;
        \begin{itemize}
            \item Simplicity. Generally, your compiler tools will do a lot of work on AST,
            and the fewer different cases you have to worry about, the better. For
            For example, if the processing of overloaded functions and operators is the same (which it is to some degree in C++), you may want to have
            only one AST node type to cover both. Having a lot of unnecessary nodes
            in the tree can be annoying as well, and may make processing slower.
            \item Good correspondence with the constructs of the language.
            \item Availability of information during processing. Some information that
            can be computed from the tree (such as type information) and might be encoded directly in the tree (at least at later stages) for easy processing.
            \item Being end-user friendly or familiar to most programmers isn't an important consideration - the abstract syntax may be radically different from
            the surface concrete syntax if that helps the compiler writer.
            \subsubsection{Sum of products}
            You may have encountered the term \gls{sum of prod}, lets's quickly run over why we use the terms \textit{sum} and \textit{product} 
            to describe the data types, and show some examples in both Haskell and Java.
            \subsubsection*{Haskell}
            \begin{lstlisting}[language=Haskell]
data SomeType = A Bool Bool Bool
                | B Bool
                | C
            \end{lstlisting}
            Here the type SomeType has 3 constructors, A, B, and C, where A takes 3 parameters, B takes one, and C zero. 
            The type of SomeType could be expressed algebraically as
            \begin{align*}
                \underbrace{(\text{Bool} \times \text{Bool} \times \text{Bool})}_{A} + \underbrace{\text{Bool}}_{B} + \underbrace{1}_{C}
            \end{align*}
            The Bool type can take on 2 different values
            ( False and True ), so the constructor
            can construct $2 * 2 * 2 = 8$ different values, since there are 8 different combinations you can
            make from 3 booleans (e.g. A True False False is one example). The constructor B can
            produce 2 different values, and C can only produce one (not zero!).
            Thus, the total numbers of values of type SomeType is $8 + 2 + 1 = 11$, as the data type
            is the sum of the three products we've just described.\\
            In short, a sum type denotes “one of” its constituent types (if your function takes SomeType
            as input, it will get either (A b1 b2 b3 ), (B b1) or C for some boolean values b1 . . .), while
            a product type denotes “all of” its constituent types (e.g. a value constructed with A will
            have all 3 booleans present.) Another way to express a product type in Haskell is with tuples,
            e.g.
            \begin{lstlisting}[language=Haskell]
type MyTriple = (Bool, Int , Char)
            \end{lstlisting}
            \subsubsection*{Java}
            In Java, we can model the same kind of data types using classes and inheritance. The instance
            variables of a class determine a “product” type, e.g.
            \begin{lstlisting}[language=Java]
class SomeClass {
    boolean a;
    boolean b;
    boolean c;
}
            \end{lstlisting}
        Similar to the constructor for A in the previous section, there are $2 * 2 * 2 = 8$ different
values an object of class SomeClass can have. Add another boolean, and you get 16 different
values. Technically, a variable of type SomeClass can take one 8 + 1 different values, since
null is also a valid value for all object variables in Java, but sometimes we ignore this fact
and tell the users of our functions to kindly not pass in null as an argument where we expect
an actual object.
Now, to get sum types, we might use class hierarchies in Java:

\begin{lstlisting}[language=Java]
interface SomeType {}

class A implements SomeType {
    boolean a;
    boolean b;
    boolean c ;
}
class B implements SomeType {
    boolean a;
}
class C implements SomeType {
}

\end{lstlisting}
    Now, an object of type SomeType can (ignoring null) take on $8 + 2 + 1$ different values,
    just like in the Haskell example above.


        \end{itemize}
        \subsection{Static Analysis}
            \subsubsection{Type Checking}
            Programming language typing always falls into one or two categories; static, and dynamic typing. Static typing is when type checking happens at compile-time, 
            while dynamic typing happens during runtime. We will focus on static typing. 
            Statically typed languages typically have these properties.
            \begin{itemize}
                \item Variables and data structures must be declared before use.
                \item Variables and data structure fields can only hold values of the declared type.
                \item Operations(i.e functions, procedures, methods) and types must be declared.
                \item Declaring the exact types of variables and operations isn't always needed. Some languages use type inference (Discussed in 3.2.3).
            \end{itemize}

            What is a type checker and how does it work? 
            A type checker is a meta-program that checks that verifies that the type of some construct(lists, expressions, etc.) matches what's expected of it.
            For example, a type checker will check that the Plus \gls{expression} takes two Integers. This lets the type checker discover and report certain errors before the program runs. 
            To do this a type checker needs to know;
            \begin{itemize}
                \item How the language should look.
                \item The language types.
                \item Rules for assigning types to the constructs.
            \end{itemize}
            \subsubsection*{Let's do a practical example}
            Here's the abstract syntax for our Simple Typed Language or STL for short.
            \begin{lstlisting}[language=Haskell]
    type VarDecl = [(String, ExprType)]

    data Expr =
        Var String
        | I Int
        | B Bool
        | BinOp Op Expr Expr
        | UnOp Op Expr
        | Choice Expr Expr Expr deriving (Show, Eq)
                
    data Op = Plus | Mult | Or | And | Not | Eq deriving (Show, Eq
            \end{lstlisting}
            Before we can continue we need to define the types that are allowed. We decide to use two types, Integers, and Booleans.
            \begin{lstlisting}[language=Haskell]
    data ExprType = Integer | Boolean deriving (Show, Eq)
            \end{lstlisting}
            Now to the meat of the exercise, the type checker itself.
            The usual way to do this in Haskell(and most other languages I've experienced) is to define a series of recursive functions, one for each
            expression/operand.
            \begin{lstlisting}[language=Haskell]
                typeCheck :: VarDecl->Expr->ExprType
            \end{lstlisting}
            I have found it easiest, to begin with, the cases that form the basic "building blocks" of a language, the literals. It is easy to know the type of Int Literals(Integer obviously), and Bool Literals(Boolean).
            \begin{lstlisting}[language=Haskell]
    typeCheck _ (I _) = Integer
    typeCheck _ (B _) = Boolean
            \end{lstlisting} 
            After these "base" cases we add a case for Vars, this is slightly more complicated since the type is dependent on the list of var declarations, 
            so we need to check the VarDecl list.
            \begin{lstlisting}[language=Haskell]
    typeCheck vars (Var name) = case lookup name vars of
                Just tp -> tp
                Nothing -> error $ "No variable could be found named "++name
            \end{lstlisting}
            We use a pretty clever Haskell expression called "case". This lets us pattern match the result of the function call in lookup. 
            I strongly recommend that you all learn how to use these since they are extremely useful and I'll be using them liberally during this example.\\
            We now move on to the ops. 
            \begin{lstlisting}[language=Haskell]
typeCheck vars (UnOp Not expr) = 
                case typeCheck vars expr of
                        Boolean-> Boolean 
                        _-> error "Argument not boolean"

typeCheck vars (BinOp Plus left right) = 
                case (typeCheck vars left, typeCheck vars right) of
                        (Integer, Integer)->Integer
                        _->error "One of the args is not an Integer"

typeCheck vars (BinOp Mult left right) = 
                case (typeCheck vars left, typeCheck vars right) of
                        (Integer, Integer)->Integer
                        _->error "One of the args is not an Integer"

typeCheck vars (BinOp Or left right) = 
                case (typeCheck vars left, typeCheck vars right) of
                        (Boolean, Boolean)->Boolean
                        _->error "One of the args is not a Boolean"

typeCheck vars (BinOp And left right) = 
                case (typeCheck vars left, typeCheck vars right) of
                        (Boolean, Boolean)->Boolean
                        _->error "One of the args is not a Boolean"
typeCheck vars (BinOp Eq left right) = Boolean
            \end{lstlisting}
            These all more-or-less follow the same pattern \footnote{Maya made me use fancier words, apparently "pretty similar" isn't good enough.}.
            They all check that the arguments are of the correct type and return the operand type if so, if not they raise an error.
            Eq is the odd one out since it returns a boolean no matter what since it checks if two expressions are the same.\\
            The last case is the most complex. Choice tests a boolean condition and returns one of the two expressions depending on the value.
            The problem is that we don't know which branch will return. We have therefore decided both branches need to be of the same type.
            \begin{lstlisting}[language=Haskell]
typeCheck vars (Choice test left right) = 
        case typeCheck vars test of
                Boolean | l == r -> r
                        | otherwise -> error "Args did not match"
                _-> error "Test condition is not a Boolean"
        where 
        l = typeCheck vars left
        r = typeCheck vars right
            \end{lstlisting}
            Here we begin by checking that the test evals to a boolean type, if so we check that both branches have the same type and return that type.
            Note that even though an evaluator would only evaluate one of the two branches, we still type check them both.
            \subsubsection{Wellformedness}
                For a program to be \gls{wellformed}it needs to satisfy all the constraints(kinda like rules) on it. This means that the program follows all the rules 
                for it like; 
                \begin{itemize}
                    \item The program conforms to the AST.
                    \item It is typed correctly
                    \item All procedure calls/declarations are wellformed.
                    \item and much more.
                \end{itemize}
                To check if a program is wellformed we usually implement a so-called constraint checker. These are pretty much just unit tests.
                The recipe for a constrain checker is as follows;
                \begin{itemize}
                    \item \textbf{Negative test cases} - Designate one negative test case for each constraint that should be checked. 
                    Ideally, each such test case should violate only one constraint.
                    \item \textbf{Reporting} - Choose an approach to "reporting". The result of constraint violation may be communicated either through
                                            a boolean value, as a list of errors, or by throwing an exception.
                    \item \textbf{Modularity} - Implement each constraint in a separate function, thereby allowing modularity and testing.
                    \item \textbf{Testing} - The constraint violations must be correctly detected for the negative test cases. The positive test case must pass.
                \end{itemize}
            \subsubsection{Type inference}
            Some languages don't specify the type of expressions and therefore need to be inferred. See example.
            Some of the expressions in the language below have different types depending on the input. We, therefore, need to create a type inferred(real word?).
            \begin{lstlisting}[language=Haskell]
type ShittyEnviroment = [(String, Expr)]

data Expr = 
        Var String
        | IntLit Integer
        | StrLit Str
        | Plus Expr Expr
        | Mult Expr Expr
            \end{lstlisting}
            We decide on some rules for those pesky ambiguous expressions. A string plus another string produces a string, an integer, and a string also produces a string.
            We also decide that it should be possible to multiply a string with a number. Using these rules we start making our type inferior.
            Like previously we start with the literals and vars.
            \begin{lstlisting}[language=Haskell]
inferType _ (IntLit _) = Integer
inferType _ (StrLit _) = String
inferType env (Var name) = case lookup name env of
                        Just expr -> inferType env expr
                        _ -> error "Can't find var"
            \end{lstlisting}
            We then implement Plus and Mult expressions based on the rules.
            \begin{lstlisting}[language=Haskell]
inferType env (Plus left right) = case (inferType env left, inferType env right) of
                                (Integer, Integer)->Integer
                                (_,_)->String --If it is not two ints it's a string!

ingerType env (Mult left right) = case (inferType env left, inferType env right) of
                                (String, Integer)-> String
                                (Integer, Integer)->Integer
                                _->error "invalid arguments"
            \end{lstlisting}
            We decide on the return type by looking at the argument types.
    \subsection{Q\&A}

    \section{State \& Scoping}
        \subsection{State}
        \subsubsection{Store}
        When running a program we often want to remember intermediary values or have variables. For us to do this we need some way of storing values.
        To do this we create a \gls{store}. A store is very simply an array, usually containing either bytes or ValueTypes. 
        The store is the program's memory(this is true in C). To access a value located at $i$ in our store we simply access the value at array index $i$.
        Worth noting that the store is what we call the program heap and by tradition, it grows upwards. This means that new entries are stored at the highest available index in our store.
        This is because the stack grows upwards and gives us the best possible use of the memory.
        \subsubsection{Variables and Enviroment}
        We now have a place to store stuff, but how do we know where it is in the store. This is where \glspl{env} comes into play.\\
        An environment is a map between variables and their location in the store. 
        \subsection{Scoping}
        A \gls{scope} is a collection of identifier bindings - i.e, what is captured by the environment at some point in the code or in time.\\
        The scope of a declaration includes all the points in the code where that declaration exists(binding). In lexical scoping the
        scope of a declaration is usually either included in the declaration constructor or extends to the nearest scoping container(think. Haskells let, where | Java curly-braces).
        In dynamic scoping, the scope of a declaration is determined at runtime and lasts until the program exits the scoping construct.
        Bindings can also be shadowed by declaring a new variable with the same name as an in-scope variable. 
        The results in the outer variables passing out of scope while the inner variable is in-scope.
        The shadowed variable still exists, but it's not accessible until the shadowing variable passes out of scope.\\
        Scoping is especially important when it comes to procedures and they usually have their own environments.\\
        \textit{Speaking of procedures\dots}
        \subsection{Q\&A}

    \section{Procedures}
    A procedure is a programming construct that abstracts away the implementation of an algorithm. Instead of writing 20 lines of code 
    every time you want to run an algorithm we simply call a procedure with those 20 lines of code inside it. This increases the readability 
    of our code and makes it easier to write it. 
    \subsection{Procedure Declarations}
    In this course, we have a simplified procedure declaration inspired by PASCAL. Our procedure declarations have a list of 
    parameters(variables that pass in/out of the procedure), another list of local variables(variables that are used within), 
    and a single statement(the algorithm itself).
    If we were to translate this abstract syntax into something more concrete then a typical procedure might look something like this.
    \begin{lstlisting}
procedure p ( upd a: integer , out b: integer ; obs c : boolean ) ;
    var x: integer ;
        y: boolean ;
        z : boolean ;
    begin
        if c then begin b := a ; x := b ; a := x end;
        y := not c ;
        c := y or c ;
    end;
    \end{lstlisting}
    The abstract syntax would be something like this in Haskell;
    \begin{lstlisting}[language=Haskell]
-- | Procedure declaration
data ProcedureDeclaration = Proc
    String -- Name of the procedure
    [Parameter] -- Parameter list
    [VarDecl] -- Local variables
    Stmt -- Statement part
    deriving (Show, Eq, Read)

-- | Procedure parameters: mode and variable declaration
type Parameter = (Mode, VarDecl)

-- | Parameter modes: observe, update, output
data Mode = Obs | Upd | Out 
            deriving (Show, Eq, Read)

-- | Variable declaration: variable name and its typ
type VarDecl = (Var,Type)

    \end{lstlisting}
    \subsubsection{Parameters \& Local Variables}
    A \gls{local var} is a variable that only exits in a specific part of the program. We're going to talk about them in the context of procedures.
    These local variables are declared or defined and only used within that procedure.
    Since the procedure exists in a vacuum, the local variable is allowed to have the same name as variables in other parts of the program. 
    Since the variable is stored at a different store location any change to the local variable wouldn't change the variable outside.

    \gls{param} are the variables the procedure uses to communicate to the outside world. They are variables, but with one key difference;
    They also specify \textit{how} it communicates with the outside world, the parameters specify if the variable is an output variable(write), observed only(read), or updatable(is this a word?)(read and write).\\
    Another way of looking at parameters is that parameters are local variables that are initialized with the passed arguments at invocation time.
    \subsubsection{Performing a Procedure}
        When performing a function we can usually assume that the parameters have already been added to the environment and initialized and that except for those parameters we have a clean
        environment to work with.
        The performing of a procedure consists of the following steps;
        \begin{enumerate}
            \item The first is that we need to somehow remember the current environment or delete all the local variables when we're done.
                The best way is to get the current \gls{stackframe}.
            \item We then add all the local variables.
            \item We then execute the procedure statements.
            \item We then reset the environment back to how it looked before using the saved stack frame from (1). This ensures that all local variables
                that shouldn't exist outside the procedure are removed. 
        \end{enumerate}
    \subsection{Executing a Procedure Call}
        We now know how to actually perform a procedure, but for us to even be able to perform a procedure in the first place we need to prepare the environment and program for it.
        There are two main ways of making the program ready to perform a procedure. These are dependent on what type of \gls{argument} passing we are doing.
    \subsection{Arg Passing}%Difference between args and params
        Argument passing fall into one of two camps.
        \subsubsection{Copy Semantics}
            The first is \gls{copy sem}. 
            In copy semantics, all the arguments are passed to the parameters by copying the value into them.
            With copy semantics, the procedure arguments are first evaluated and then bound to the parameters when they are added.
            The procedure is then performed, and the results are obtained. Those are then copied back to all the arguments that are upd or out.
            To further explain let's look at a trivial example to understand what happens in-store.
            Say we have a procedure that simply adds 5 to any number.
\begin{lstlisting}
procedure add5(obs x: integer, out res: integer)
begin
    res = x + 5;
end;

\end{lstlisting}
\newpage
            Now let's look at the main procedure
\begin{lstlisting}
procedure main()
var a: integer;
var b: integer;
begin
    a = 42;
    add5(a, b);
end;
\end{lstlisting}
    \Cref{fig:mainBeforeCall} shows us what the state looks like before we start executing the procedure call.
    \begin{figure}[!h]
        \StateDiagram{97}{97}{U,U,I 42}{b/98,a/99}
        \caption{State before the procedure call}
        \label{fig:mainBeforeCall}
    \end{figure}
    The first thing we do when entering the procedure call is to save the stackframe so we can restore the environment to its proper place once we're done. We then clear the environment, ensuring that the only variable declared by the procedure is in scope.
    We then add all the parameters(in our case we add x and res) and copy the values of a and b into x and res respectively \footnote{since res is out it isn't initialized}.
    \Cref{fig:mainBeginCallCopy} shows what the state looks like at this point.
    
    \begin{figure}[!h]
        \StateDiagram{95}{95}{U,U,I 42,U,I 42}{res/96,x/97}
        \caption{State when entering add5 (copy semantics).}
        \label{fig:mainBeginCallCopy}
    \end{figure}     

    We have now prepared the procedure so that it can be performed.
    \Cref{fig:mainEndCallCopy} shows how the state looks after its been performed.
    \begin{figure}[!h]
        \StateDiagram{95}{95}{U,I 47,I 42,U,I 42}{res/96,x/97}
        \caption{State after res = x +5 in add5 (copy semantics).}
        \label{fig:mainEndCallCopy}
    \end{figure}

    Now comes the fun part, \textit{cleaning}! Just performing the function isn't enough, if we were to just drop it here we would have an environment that 
    is drastically different than then when we started. We could just reset the stackframe to what it was when we started, but then b wouldn't get its new value. 
    What we need to do is copy the values of the out/upd params back to their argument variables. Then we can reset the stackframe making our state look like
    \cref{fig:mainBeforeCall}. As you can see the free pointer now points to where x used to point and b now has the same value as a res. That means that the next time we add something to the environment it will
    overwrite those values in the store.     
        
    \begin{figure}[!h]
        \StateDiagram{95}{97}{U,I 47,I 42,I 47,I 42}{b/98,a/99}
        \caption{State after add5(a, b) in main (copy semantics).}
        \label{fig:mainAfterCall}
    \end{figure}

        \subsubsection{Reference Semantics}
            Now, what if \verb+add5+ could instead write \emph{directly} to b
            in \verb+main+ instead of having to allocate its own copy which is later copied?
            This is the main advantage of \gls{ref sem}. Lets reuse the example from before.
            As you can see the from \cref{fig:mainBeforeCall2} the state before we enter the procedure call is the same as during copy semantics.
            \begin{figure}[!h]
                \StateDiagram{97}{97}{U,U,I 42}{b/98,a/99}
                \caption{State before the procedure call}
                \label{fig:mainBeforeCall2}
            \end{figure}
            
            It is now we encounter our first change. Like before we get the stackframe before clearing the environment and adding the parameters, but here comes the change.
            instead of allocating space and copying the value of the args to the upd/out params we instead set their address to the address of the corresponding argument.
            We now get a state that looks like \cref{fig:mainBeginCallRef}. Since x is an obs parameter it is still allocated and copied to as normal.
            \begin{figure}[!h]
                \StateDiagram{96}{96}{U,I 42,U,I 42}{res/98,x/97}
                \caption{State when entering add5 (reference out).}
                \label{fig:mainBeginCallRef}
            \end{figure}

            Since res now points to the same place as b, any changes to ref will also be reflected by b. This makes cleanup much easier!
            infact all we have to do after performing the procedure is to reset the enviroment back to where it was and call it a day makng the state look like
            \cref{fig:mainEndCallRef}.
        \begin{figure}[!h]
            \StateDiagram{96}{96}{U,I 42,I 47,I 42}{res/98,x/97}
            \caption{State after res = x+5 in add5 (reference out).}
            \label{fig:mainEndCallRef}
        \end{figure}

    \subsection{Q\&A}

    \section{Generics}
        To explain generics we'll demonstrate an example where generics are useful. Say we wanted to test if some algebraic type was a ring.
        It would be inefficient to create a unique test for each algebraic type we wanted to check. This is where generics are useful.
        A generic mechanism would let us write only one of these tests. 
        Any generic mechanism needs three components, a
        Lets start with showing the ASR and datastructures we're going to be using.
        \begin{lstlisting}[language=Haskell]
type Enviroment valuedomain= [(String, ExprAST valuedomain)]
data ExprAST valuedomain
    = Lit valuedomain TypeName
    | Fun FunName [ExprAST valuedomain]
    | Var VarName
    | Assert (ExprAST valuedomain)
    deriving (Eq, Read, Show)
        \end{lstlisting}
        Now for an algebraic structure to be a ring, certain properties need to hold. We, therefore, add some tests to check for these\footnote{The implementation is found in Appendix B}.
        \begin{lstlisting}[language=Haskell]
--Tests that carrier set is an abelian group
testAssoc :: ExprAST valuedomain
testAddComm ::ExprAST valuedomain
testAddID ::ExprAST valuedomain
testAddInv ::ExprAST valuedomain

--Tests that carrier set is a monoid under multiplication
testMultAssoc::ExprAST valuedomain
testMultID ::ExprAST valuedomain

--Test that multiplication is distributive
testLDist :: ExprAST valuedomain
testRDist :: ExprAST valuedomain

--List of tests
tests = [testAssoc, testAddComm, testAddID, testAddInv, testMultAssoc, testMultID, testLDist, testRDist]

        \end{lstlisting}
        \subsection{Signature}
        To represent the different types(of algebraic structures) we introduce the concept of a signature. 
        \begin{lstlisting}[language=Haskell]
type Signature = ([TypeDeclaration],[FunDeclaration])
        \end{lstlisting}
        A signature is a data structure that contains a list of types, and the operations(functions, etc.) that operate on those types.
        Signatures have other applications other than just generics. Defining functions in signatures means that if we want to expand our language with a new intrinsic function we won't
        have to change the code anywhere else than in the signature and the accompanying semantic implementation.\\
        Let's take a look at the signature for the Integers.
        \begin{lstlisting}[language=Haskell]
int ::Signature
int = (["Int", "Bool"],[
                ("Add", ["Int","Int"],"Int"),
                ("Mult", ["Int","Int"], "Int"),
                ("Neg", ["Int"], "Int"),
                ("Eq", ["Int", "Int"], "Bool"),
                ("LEq", ["Int", "Int"], "Bool")
        ])
        \end{lstlisting}
        As you see the integers have 5 functions, the first two are binary operators that are needed for a ring, the \verb+neg+ function is to represent the negative numbers.
        \verb+Eq+ and \verb+LEq+ are needed so we can compare the numbers and check transitivity.
        Since we have multiple types defined in our signatures we need to tell Haskell about them, this is why we have \verb+valuedomain+ everywhere.
        But since we have multiple types in the signature we need to create a custom data structure to represent it. \\
        \begin{lstlisting}[language=Haskell]
-- | The semantic domain for the rings sigs
data Ring = BoolDom Bool
        | NatDom Int Bool
        | IntDom Int Bool
        | RealDom Float Bool
        deriving(Show, Eq, Ord)
        \end{lstlisting}
        Now that we have defined what functions these algebraic structures have we need some way of defining their semantics, i.e what they do.
        To do this we create a \verb+FunctionModel+ for each function implemented by a signature.
        \begin{lstlisting}[language=Haskell]
type FunModel valuedomain = FunName -> [valuedomain] -> valuedomain
        \end{lstlisting}
        Here is the semantics of our IntegerRing
        \begin{lstlisting}[language=Haskell]
--Semantics of a ring(when carrier set is a int)
intRingModel ::FunModel Ring
intRingModel "Add" arg =  IntDom (foldl (+) 0 (map unpackInt arg)) False
intRingModel "Mult" arg = IntDom (foldl (*) 1 (map unpackInt arg)) False
intRingModel "Neg" [x] = IntDom (negate (unpackInt x)) $ getBool x
intRingModel "Eq" [x, y] = BoolDom (x==y)
intRingModel "LEq" [x, y] = BoolDom (x<=y)
        \end{lstlisting}
        We now have all we need to generify the evaluator.
        \begin{lstlisting}[language=Haskell]
-- | Eval
eval ::FunModel Ring->Signature->Enviroment Ring->ExprAST Ring->Ring
eval mod sig env (Lit i _) = i
eval mod sig env (Var name) = case lookup name env of
                            Just x -> eval mod sig env x
                            Nothing -> error "No such var in env"

eval mod (types, funcs) env (Fun name args) | elem name (map (\(x,y,z)->x) funcs) =mod name (map (eval mod (types, funcs) env) args)
                                            | otherwise= error "Can't find function in signature"

eval mod (sig) env (Assert expr) = let res@(BoolDom b) = eval mod sig env expr in if b then res else error $"Assert did not hold"++ (show expr)

        \end{lstlisting}
        Instead of the evaluator handling the evaluation of functions we pass it to the function model letting it deal with it according to the implementation defined there.
        \subsection{Property based assertions}
        In the example above we used asserts to check if an algebraic structure was a ring or not. Having asserts in a language can be extremely useful since it lets us check if certain conditions hold, and if they don't, to stop execution.
        A program crash is often the best course of action to prevent illegal states. 
        \subsection{Q\&A}

    \section{Language Standards}
        \subsection{Software Engineering Implications of Languages}
        The choice of language can have a massive impact on software development. The move to higher abstraction languages has caused a corresponding increase in efficiency
        since the developer can focus all their efforts on what the program should accomplish instead of having to work on the details. 
       \gls{SLE} has many applications within software engineering like;
        \begin{itemize}
            \item Design, implementation, and usage of DSLs that are tailor-made for a specific problem or technical domains, like, UI, web services, configurations, testing,
                    date exchange, interoperability, deployment, and distribution.
            \item Software reverse engineering and re-engineering in many forms, for example, analysis of projects regarding their dependence on open source software,
                    integration of systems, and migration of systems constrained by legislation or technology.
            \item Data extraction in the context of data mining, information retrieval, machine learning, big data analytics, social science, digital forensics, and AI, 
                    with diverse input, artifacts to be parsed interchange formats to conform to.
        \end{itemize}
        Software languages can also impact the security of any product. If a language has a fundamental flaw or creates unknown pitfalls that might be hard to spot due to
        the inherent design of the language could then be exploited by malicious actors. 
        \subsection{Reading spesifications}
        \subsubsection{Backus-Naur Form}
            \gls{BNF} is a format for describing context-free grammar. We use it to describe the syntax of languages. Although it is probably not part of the curriculum
            it's still important to know since most language specification describes their languages using some variation of BNF most common of which is the Extended Backus-Naur form(EBNF).
            An EBNF consists of terminal symbols and non-terminal production rules which are the restrictions governing how terminal symbols can be combined into a legal sequence. 
            Examples of terminal symbols include alphanumeric characters, punctuation marks, and whitespace characters. These constructions often end up looking like Haskell data structures,
            this should hopefully help you understand them.
            \begin{lstlisting}[language=XML]
<symbol> ::= __expression__
            \end{lstlisting}
            Here is the complete PASCAL-like language that only allows assignments. in its EBNF form.
            \begin{lstlisting}[language=XML]
(* a simple program syntax in EBNF - Wikipedia *)
program = 'PROGRAM', white space, identifier, white space, 
            'BEGIN', white space, 
            { assignment, ";", white space }, 
            'END.' ;
identifier = alphabetic character, { alphabetic character | digit } ;
number = [ "-" ], digit, { digit } ;
string = '"' , { all characters - '"' }, '"' ;
assignment = identifier , ":=" , ( number | identifier | string ) ;
alphabetic character = "A" | "B" | "C" | "D" | "E" | "F" | "G"
                        | "H" | "I" | "J" | "K" | "L" | "M" | "N"
                        | "O" | "P" | "Q" | "R" | "S" | "T" | "U"
                        | "V" | "W" | "X" | "Y" | "Z" ;
digit = "0" | "1" | "2" | "3" | "4" | "5" | "6" | "7" | "8" | "9" ;
white space = ? white space characters ? ;
all characters = ? all visible characters ? ;
            \end{lstlisting}
            EBNF also includes regex like syntax for expressing repetition(*,+), optionality(?,[]), and alternatives(|).
            If the above looks confusing or unintuitive I wouldn't worry. The syntax of EBNF is fairly intuitive and most standards have fairly little of it at once.
        \subsection{Q\&A}

    \section{Misc.}
        \subsection{Datastructures in memory}
        \subsubsection{Arrays}
            Say we have an array of $n$ elements of type $T$, and the type $T$ takes up 4 bytes of memory. Lets also assume that the header is 8 bytes
            The array would then look like this in the store.
            
        \subsubsection{Records}
        \subsubsection{Pointers \& linked lists}
        \subsection{Q\&A}

    \appendix
    \section{Glossary}
    \printnoidxglossaries

    \section{Code Examples}
    \subsection{Type Check Example}
    \begin{lstlisting}[language=Haskell]
module TypeCheckExample where
type VarDecl = [(String, ExprType)]

data Expr = 
        Var String
        | I Int
        | B Bool
        | BinOp Op Expr Expr
        | UnOp Op Expr
        | Choice Expr Expr Expr deriving (Show, Eq)
            
data Op = Plus | Mult | Or | And | Not | Eq deriving (Show, Eq)

data ExprType = Integer | Boolean deriving (Show, Eq)

typeCheck ::VarDecl->Expr->ExprType
typeCheck _ (I _) = Integer
typeCheck _ (B _) = Boolean

typeCheck vars (Var name) = 
                case lookup name vars of
                        Just tp -> tp
                        Nothing -> error $ "No variable could be found named "++name

typeCheck vars (UnOp Not expr) = 
                        case typeCheck vars expr of
                                Boolean-> Boolean 
                                _-> error "Argument not boolean"

typeCheck vars (BinOp Plus left right) = 
                        case (typeCheck vars left, typeCheck vars right) of
                                (Integer, Integer)->Integer
                                _->error "One of the args is not an Integer"

typeCheck vars (BinOp Mult left right) = 
                        case (typeCheck vars left, typeCheck vars right) of
                                (Integer, Integer)->Integer
                                _->error "One of the args is not an Integer"

typeCheck vars (BinOp Or left right) = 
                        case (typeCheck vars left, typeCheck vars right) of
                                (Boolean, Boolean)->Boolean
                                _->error "One of the args is not a Boolean"

typeCheck vars (BinOp And left right) = 
                        case (typeCheck vars left, typeCheck vars right) of
                                (Boolean, Boolean)->Boolean
                                _->error "One of the args is not a Boolean"

typeCheck vars (BinOp Eq left right) = Boolean

typeCheck vars (Choice test left right) = 
                        case typeCheck vars test of
                                Boolean | l == r -> r
                                        | otherwise -> error "Args did not match"
                                _-> error "Test condition is not a Boolean"
                        where
                                l = typeCheck vars left
                                r = typeCheck vars right
    \end{lstlisting}
    \subsection{Type Inference Example}
    \begin{lstlisting}[language=Haskell]
module TypeCheckExample where

type ShittyEnviroment = [(String, Expr)]

data Expr = 
        Var String
        | IntLit Integer
        | StrLit Str
        | Plus Expr Expr
        | Mult Expr Expr

--New Rules
--Int + string = string
--String+String = string
--String * int = String

data ExprType = Integer | String deriving (Show, Eq)

inferType:: ShittyEnviroment->Expr->ExprType
inferType _ (IntLit _) = Integer
inferType _ (StrLit _) = String

inferType env (Var name) = case lookup name env of
                        Just expr -> inferType env expr
                        _ -> error "Can't find var"

inferType env (Plus left right) = case (inferType env left, inferType env right) of
                                (Integer, Integer)->Integer
                                (_,_)->String --If it's not two ints it's a string!

ingerType env (Mult left right) = case (inferType env left, inferType env right) of
                                (String, Integer)-> String
                                (Integer, Integer)->Integer
                                _->error "invalid arguments"
    \end{lstlisting}
    \subsection{Generic Example}
    \begin{lstlisting}[language=Haskell]
module GenericsExample where

type TypeName = String
type FunName = String
type VarName = String

type Enviroment valuedomain= [(String, ExprAST valuedomain)]
--Defines semantics of a signature
type FunModel valuedomain = FunName -> [valuedomain] -> valuedomain

type FunDeclaration = (FunName, [TypeName], TypeName)
type TypeDeclaration = TypeName

--Defines what functions types exists
type Signature = ([TypeDeclaration],[FunDeclaration])

--DS??
data ExprAST valuedomain
    = Lit valuedomain TypeName
    | Fun FunName [ExprAST valuedomain]
    | Var VarName
    | Assert (ExprAST valuedomain)
    deriving (Eq, Read, Show)


--DSL Describing Ints
int ::Signature
int = (["Int", "Bool"],[
                ("Add", ["Int","Int"],"Int"),
                ("Mult", ["Int","Int"], "Int"),
                ("Neg", ["Int"], "Int"),
                ("Eq", ["Int", "Int"], "Bool"),
                ("LEq", ["Int", "Int"], "Bool")
        ])

--DSL for Describing Reals
real :: Signature
real =(["Real", "Bool"],[
                ("Add", ["Real","Real"],"Real"),
                ("Mult", ["Real","Real"], "Real"),
                ("Neg", ["Real"],"Real"),
                ("Eq", ["Real", "Real"], "Bool"),
                ("LEq", ["Real", "Real"], "Bool")
    ])
--Natural nums
nat :: Signature 
nat =(["Nat","Bool"],[
                ("Add", ["Nat","Nat"],"Nat"),
                ("Mult", ["Nat","Nat"], "Nat"),
                ("Eq", ["Nat", "Nat"], "Bool"),
                ("LEq", ["Nat", "Nat"], "Bool")
    ])
--Natural nums
bool :: Signature 
bool =(["Bool"],[
            ("Add", ["Bool","Bool"],"Bool"),
            ("Mult", ["Bool","Bool"], "Bool"),
            ("Eq", ["Bool", "Bool"], "Bool"),
            ("LEq", ["Bool", "Bool"], "Bool")
        ])

-- | The semantic domain for the rings sigs
data Ring = BoolDom Bool
            | NatDom Int Bool
            | IntDom Int Bool
            | RealDom Float Bool
            deriving(Show, Eq, Ord)

unpackInt :: Ring->Int
unpackInt (IntDom i _) = i
unpackNat :: Ring->Int
unpackNat (NatDom n _) = n
unpackReal :: Ring->Float
unpackReal (RealDom x _) = x
unpackBool :: Ring->Bool
unpackBool (BoolDom q) = q

getBool ::Ring->Bool
getBool (IntDom _ q) = q
getBool (NatDom _ q) = q
getBool (RealDom _ q) = q
getBool (BoolDom q) = q


-- | Semantics 

--Semantics of a ring(when carrier set is a int)
intRingModel ::FunModel Ring
intRingModel "Add" arg =  IntDom (foldl (+) 0 (map unpackInt arg)) False
intRingModel "Mult" arg = IntDom (foldl (*) 1 (map unpackInt arg)) False
intRingModel "Neg" [x] = IntDom (negate (unpackInt x)) $ getBool x
intRingModel "Eq" [x, y] = BoolDom (x==y)
intRingModel "LEq" [x, y] = BoolDom (x<=y)

--Semantics of a ring(when carrier set is a nat)
natRingModel ::FunModel Ring
natRingModel "Add" arg = NatDom  (foldl (+) 0 (map unpackNat arg)) False
natRingModel "Mult" arg = NatDom (foldl (*) 1 (map unpackNat arg)) False
natRingModel "Neg" [x] = NatDom (negate (unpackNat x)) $ getBool x
natRingModel "Eq" [x, y] = BoolDom (x==y)
natRingModel "LEq" [x, y] =BoolDom (x<=y)

--Semantics of a ring(when carrier set is a real)
realRingModel ::FunModel Ring
realRingModel "Add" arg =  RealDom (foldl (+) 0 ((map unpackReal arg))) False
realRingModel "Mult" arg = RealDom (foldl (*) 1 ((map unpackReal arg))) False
realRingModel "Neg" [x] = RealDom (negate (unpackReal x))  $ getBool x
realRingModel "Eq" [x, y] = BoolDom (x==y)
realRingModel "LEq" [x, y] =BoolDom (x<=y)

--Semantics of a ring(when carrier set is bool)
boolRingModel ::FunModel Ring
boolRingModel "Add" arg =  BoolDom $ foldl (/=) False (map unpackBool arg)
boolRingModel "Mult" arg = BoolDom $ foldl (&&) True (map unpackBool arg)
boolRingModel "Neg" [x] = BoolDom (not (unpackBool x))
boolRingModel "Eq" [x, y] = BoolDom (x==y)
boolRingModel "LEq" [x, y] =BoolDom (x<=y)


-- | Eval
eval ::FunModel Ring->Signature->Enviroment Ring->ExprAST Ring->Ring
eval mod sig env (Lit i _) = i
eval mod sig env (Var name) = case lookup name env of
                            Just x -> eval mod sig env x
                            Nothing -> error "No such var in env"

eval mod (types, funcs) env (Fun name args) | elem name (map (\(x,y,z)->x) funcs) =mod name (map (eval mod (types, funcs) env) args)
                                            | otherwise= error "Can't find function in signature"

eval mod (sig) env (Assert expr) = let res@(BoolDom b) = eval mod sig env expr in if b then res else error $"Assert did not hold"++ (show expr)


-- | Tests that carrier set is abelian group
testAssoc :: ExprAST valuedomain--We use Integer to represent the set 
testAssoc = 
    Assert (
        Fun "Eq" [
            Fun "Add" [
                Fun "Add" [
                    Var"a",
                    Var"b"
                    ], 
                Var"c"
            ],
            Fun "Add" [
                Var"a", 
                Fun "Add" 
                [
                    Var"b", 
                    Var"c"
                ]
            ]
        ]
    )

testAddComm ::ExprAST valuedomain
testAddComm = 
    Assert (
        Fun "Eq" [
            Fun "Add" [
                Var "a",
                Var "b"
            ],
            Fun "Add" [
                Var "b",
                Var"a"
            ]
        ]
    )

testAddID ::ExprAST valuedomain
testAddID = 
    Assert (
        Fun "Eq" [
            Fun "Add" [
                Var "a",
                Var "addConst"],
            Var"a"
        ]
    )

testAddInv ::ExprAST valuedomain
testAddInv = 
    Assert (
        Fun "Eq" [
            Fun "Add" [
                Var"a",
                Fun "Neg" [
                    Var"a"
                ]
            ],
            Var "addConst"
        ]
    )

--Tests that carrier set is a monoid under multiplication
testMultAssoc::ExprAST valuedomain
testMultAssoc = 
    Assert (
        Fun "Eq" [
            Fun "Mult" [
                Fun "Mult" [
                    Var"a",
                    Var"b"
                ],
                Var "c"
            ],
            Fun "Mult" [
                Var"a", 
                Fun "Mult" [
                    Var"b",
                    Var"c"
                ]
            ]
        ]
    )

testMultID ::ExprAST valuedomain
testMultID = 
    Assert (
        Fun "Eq" [
            Fun "Mult" [
                    Var"a", 
                    Var "multConst"
                ],
            Var "a"
        ]
    )

--Test that multiplication is distributive
testLDist :: ExprAST valuedomain
testLDist = 
    Assert (
        Fun "Eq" [
            Fun "Mult" [
                Var "a", 
                Fun "Add" [Var "b",Var "c"]
            ], 
            Fun "Add" [
                (Fun "Mult" [Var"a",Var "b"]),
                (Fun "Mult" [Var"a",Var "c"])
            ] 
        ]
    )

testRDist :: ExprAST valuedomain
testRDist =  
    Assert ( 
        Fun "Eq" [
            Fun "Mult" [
                Fun "Add" [Var "b", Var"c"], 
                Var"a"
                ], 
            Fun "Add" [
                (Fun "Mult" [Var "b",Var "a"]),
                (Fun "Mult" [Var "c",Var "a"])
                ] 
        ]
    )



tests = [testAssoc, testAddComm, testAddID, testAddInv, testMultAssoc, testMultID, testLDist, testRDist]

-- | main
test tpname (sig, model) dta = let env = [("a", Lit (dta!!0) tpname), ("b", Lit (dta!!1) tpname), ("c", Lit (dta!!2) tpname), ("multConst", Lit (dta!!3) tpname), ("addConst", Lit (dta!!2) tpname)] in
    map (eval model sig env) tests
    \end{lstlisting}
\end{document}